\newif\ifglobal

%%%%%%%%%%%%%%%%%%%%%%%
%%% Compile to view %%%
%%%%%%%%%%%%%%%%%%%%%%%

%%%%%%%%%%%%%%%%%%%%%%%%%%%%%%%%%%%%%%%%%%%%%%%%%%%
%%% Readme:                                     %%%
%%% https://www.overleaf.com/read/bwdjvfjksvjy/ %%%
%%%%%%%%%%%%%%%%%%%%%%%%%%%%%%%%%%%%%%%%%%%%%%%%%%%

% >>> M?: marks a module setting number or important notice

% >>> M4: Set {./relative-path} to external files for this module <<<
\def \modulefiles {./01-BUSARCH}
% To add an external file, use:
% (Replace '---' with actual filename)
% '\input{\modulefiles/tables/---.tex}' for tables
% '\includegraphics{\modulefiles/figures/---}' for figures
% '\subfile{\modulefiles/---}' for register files

% >>> M1: This module is in global project? <<<
% 'YES' - uncomment; 'NO' - comment out
\globaltrue

\ifglobal
    \documentclass[main\_\_EN.tex]{subfiles}
   
\else
    \documentclass[openany, 10pt]{book}

    % >>> M2: In ./00-shared/preamble-trm-module, also find
    %         settings for visibility of todo notes, watermark <<<
    \usepackage{preamble-trm-module}

    % >>> M3: Temporary labels in English,
    %         you can add your own labels there <<<
    \myexternaldocument{./00-shared/temp-labels-en}

\fi %\ifglobal


\setcounter{secnumdepth}{4}


% >>> M5: If this module needs any updates in
% './00-shared/preamble-trm-module.sty'
% (include additional package, etc.), contact your module PM
% or create the file './00-shared/changelog.tex' make the
% necessary updates yourself and document them in this file <<<

% Make text left-aligned and allow latex to break pages naturally, comment out for full justification
\raggedright
\raggedbottom


\begin{document}


% Sets language into which pre-defined bits of text are translated
\selectlanguage{english}
\setenglishformatting

% Do not delete this block
\ifglobal
% Add the following front matter if
% this module is outside of global project
\else
    \InputIfFileExists{readme}{}{}
    \clearpage
    \listoftodos
    \listofdonetodos
    \tableofcontents
    %\listoftables
    %\listoffigures
\fi


%%%%%%%%%%%%%%%%%%%%%%%%%
%%% TRM Module Begins %%%
%%%%%%%%%%%%%%%%%%%%%%%%%

\hypertarget{busarch}{}
\chapter{Bus Architecture}
\label{mod:busarch}

The \chipname{} SoC uses a multi-layer bus matrix as the core of its communication architecture, connecting multiple masters and slaves through an efficient interconnect. This architecture supports concurrent access and reduces access conflicts, so that the high-performance and low-power RISC-V cores, DMA controllers, debug interfaces, and various peripherals can operate at the same time with minimal contention. By optimizing data flow and system performance, the bus matrix keeps the system efficient even when multiple high-speed peripherals operate simultaneously.

The main system of the \chipname{} SoC consists of:

\begin{itemize}
    \item Masters:
        \begin{itemize}
            \item I-bus, D-bus, and AHB bus of high-performance (HP) 32-bit RISC-V processor 0
            \item I-bus, D-bus, and AHB bus of high-performance (HP) 32-bit RISC-V processor 1
            \item I-bus, D-bus, and AHB bus of the low-power (LP) 32-bit RISC-V processor
            \item TCM MEM MONITOR
            \item PSRAM MEM MONITOR
            \item TRACE0
            \item TRACE1
            \item GDMA-AHB-HP
            \item GDMA-AHB-LP
            \item GDMA-AXI
            \item USB OTG\_HS
            \item GMAC
            \item SD/MMC
            \item 2D-DMA
        \end{itemize}
    \item Slaves:
        \begin{itemize}
            \item HP \& LP peripherals
            \item CPU peripherals
            \item ROM
            \item LP SRAM
            \item HP SRAM
            \item External RAM
            \item External flash
        \end{itemize}
\end{itemize}

The bus architecture of \chipname{} is shown in the figure below.

\begin{figure}[H]
    \centering
    \includegraphics[width=1.0\textwidth]{\modulefiles/figures/bus_matrix.pdf}
    \begin{minipage}{0.9\textwidth}
    \small
    Note: Circles at intersections indicate that the bus in that column can access the slave in the corresponding row.\\
    All masters can only read from external flash.\\
    \end{minipage}
    \caption{\chipname{} Bus Architecture}
    \label{fig:bus-arch}
\end{figure}

The buses in Figure \ref{fig:bus-arch} \textit{\nameref{fig:bus-arch}} are described as follows:

\begin{itemize}
    \item \textbf{HP CPU I-bus}: Connects the instruction buses of high-performance (HP) processor 0 and processor 1 to the bus matrix for instruction fetch. Accessible targets include ROM, HP SRAM, external RAM, and external flash. Accesses to external memory go through the cache.
    \item \textbf{HP CPU D-bus}: Connects the data buses of high-performance (HP) processor 0 and processor 1 to the bus matrix for data loads and stores by the processor and memory access by the debugger. Accessible targets include ROM, HP SRAM, external RAM, and external flash. Accesses to external memory go through the cache.
    \item \textbf{HP CPU AHB}: Connects the AHB buses of high-performance (HP) processor 0 and processor 1 to the bus matrix for accessing data in peripherals or memory. Accessible targets include HP \& LP peripherals, CPU peripherals, LP SRAM, external RAM, and external flash. Accesses to external RAM and external flash through this bus are used for debugging. Accesses to external memory through this bus bypass the cache. The access addresses are remapped to the address range used when accessing through the cache. For details about address mapping, see Chapter \ref{mod:sysmem} \textit{\nameref{mod:sysmem}}.
    \item \textbf{LP CPU I-bus and D-bus}: Connect the instruction bus and data bus of the low-power (LP) processor to the bus matrix for instruction fetch and data access respectively. The accessible target is LP SRAM.
    \item \textbf{LP CPU AHB}: Connects the AHB bus of the low-power (LP) processor to the bus matrix for accessing data in peripherals or memory. Accessible targets include HP \& LP peripherals, CPU peripherals, HP SRAM, external RAM, and external flash.
    \item \textbf{TCM MEM MONITOR}: Connects the TCM memory monitor to the bus matrix for writing monitoring data to memory. Accessible targets include HP SRAM and external RAM.
    \item \textbf{PSRAM MEM MONITOR}: Connects the PSRAM memory monitor to the bus matrix for writing monitoring data to memory. Accessible targets include HP SRAM and external RAM.
    \item \textbf{TRACE}: Connects the TRACE0 and TRACE1 modules to the bus matrix for debug operations. The two TRACE interfaces have the same access path. Accessible targets include HP SRAM and external RAM.
    \item \textbf{GDMA}: The chip has three DMA controllers for high-speed data transfer: GDMA-AHB-HP, GDMA-AHB-LP, and GDMA-AXI.
        \begin{itemize}
            \item \textbf{GDMA-AHB-HP}: DMA controller based on the AHB bus protocol. Accessible targets include HP SRAM, external RAM, and external flash.
            \item \textbf{GDMA-AHB-LP}: DMA controller based on the AHB bus protocol. Accessible targets include LP SRAM, HP SRAM, external RAM, and external flash.
            \item \textbf{GDMA-AXI}: DMA controller based on the AXI bus protocol. Accessible targets include HP SRAM, external RAM, and external flash.
        \end{itemize}
        GDMA transfers data over these buses: peripheral-to-memory, memory-to-peripheral, and memory-to-memory. For DMA accesses to external memory, use a longer burst length to keep throughput up.
    \item \textbf{USB OTG\_HS}: The high-speed USB OTG controller, whose built-in DMA is connected to the bus matrix, for USB 2.0 high-speed data transfer. Accessible targets include HP SRAM, external RAM, and external flash.
    \item \textbf{GMAC}: The Ethernet MAC controller, whose built-in DMA is connected to the bus matrix, for network data transfer. Accessible targets include HP SRAM, external RAM, and external flash.
    \item \textbf{SD/MMC}: The SD/MMC controller, whose built-in DMA is connected to the bus matrix, for SD/MMC card read and write operations. Accessible targets include HP SRAM, external RAM, and external flash.
    \item \textbf{2D-DMA}: Connects the 2D graphics accelerator to the bus matrix for 2D graphics data transfer and processing. Accessible targets include HP SRAM, external RAM, and external flash.
\end{itemize}

\chipname{}'s cache is used only for CPU access to external memory. The CPU I-bus and D-bus access external RAM through the cache, while masters such as GDMA, USB OTG\_HS, GMAC, SD/MMC, and 2D-DMA access external RAM directly, without going through the cache. Hardware does not maintain data coherence between these paths. Therefore, software must ensure coherence when the CPU and these masters share data in external RAM.

Data can become inconsistent in two ways:

\begin{itemize}
    \item Data written by the CPU may remain in the cache and not yet be written back to external RAM. In this case, a master may read stale data. Before a master reads the data, software must write the corresponding cache contents back to external RAM.
    \item After a master writes to external RAM, the cache may still contain a stale copy of the data. In this case, the CPU may read stale data. Before the CPU reads the data, software must invalidate the corresponding cache contents. If the master writes executable code, the instruction cache must also be invalidated.
\end{itemize}

Internal memory is not accessed through the cache, and all masters can only read from external flash. Therefore, neither is subject to the issue described above.

\end{document}